\documentclass[11pt,a4paper]{article}
\usepackage[margin=25mm]{geometry}
\usepackage{amsmath,amssymb,amsthm,hyperref}
\setlength{\emergencystretch}{2em}
\newtheorem{proposition}{Proposition}
\newcommand{\gen}[1]{\langle#1\rangle}
\title{Excluding transitive-side equality under an alternating tail}
\author{Human Compression project: bounded L044 proof component}
\date{}
\begin{document}\maketitle
Let $G\leq A_n$, $n\geq12$, have a transitive string C-marking
with all maximal omissions intransitive. Suppose a unique eligible
interior $i$-edge $ab$ joins the two omission orbits $A,B$, and
$L=G_{<i}$ fixes $B$ pointwise and is transitive on $A$.
Assume the equality configuration of the accepted transitive-side
support theorem: with $h=i+1<r-1$ and $K=G_{>h}$,
\[
 |A|=2h,\qquad \operatorname{Fix}(K)=A\cup\{b\}.
\]
Put $X=B\setminus\{b\}$, and suppose $K$ is the full natural
alternating group $A_X$. Write $m=h=i+1\geq2$ and
$\tau=\rho_{i+1}$.

\begin{proposition}
The stated equality configuration is impossible.
\end{proposition}
\begin{proof}
The accepted transitive-side proof shows that $\tau$ is
fixed-point-free on $A$ and centralizes $L$. Its pairs form
an $L$-invariant system of $m$ blocks. Every distinguished
generator of $L$ has essential block image, by the actual
global-omission argument in that proof. These $m-1$ images
are independent, so Whiston's equality clause makes the
block image the full $S_m$.

Write $\lambda=\tau|_A$, $a'=a^\lambda$, and
$H=G_{<i-1}$. The generators of $H$ fix $a$ by unique crossing
and fix $a'$ by commutation with $\tau$. Their $m-2$
independent block images fix the pair $\{a,a'\}$ and generate
the full $S_{m-1}$ on the others, again by Whiston. For $m=2$
this is the trivial $S_1$ statement. No classification of a
maximal symmetric string and no projected IP are used.

Since the nonempty involutory family generating $K=A_X$
cannot generate $A_3$ or $A_4$, one has $|X|\geq5$.
The crossing generator $s=\rho_i$ centralizes $K$, preserves
its unique nontrivial orbit $X$, and fixes $X$ pointwise
because the natural alternating centralizer is trivial.
Consequently
\[
                      s=(a\ b)u,
\]
where $u$ is supported on $A\setminus\{a\}$ and centralizes
$H$. For $m\geq3$ the only $H$-fixed points in $A$ are
$a,a'$, so $u$ fixes $a'$ as well.

First suppose $m$ is even. The restriction $\tau|_B$ is even
and moves $b$ into $X$. The natural alternating groups on
$B\setminus\{b\}$ and $B\setminus\{b^\tau\}$ generate
$A_B$, by the elementary overlapping-support three-cycle
argument. Thus $R=G_{>i}=\gen{\tau,K}$ is a subgroup of
$\gen\lambda\times A_B$ projecting onto both factors.
Perfection of $A_B$ implies $\lambda\in R$: otherwise
the projection to $A_B$ would identify $R$ with a graph of
a nontrivial homomorphism $A_B\to C_2$.
Since $L$ and $R$ have disjoint index sets, IP gives
$\lambda\notin L$.

Identify the pair stabilizer with $\mathbb F_2^m\rtimes S_m$.
Its base vectors record the within-pair flips; block
permutations are even on the $2m$ points. Hence the kernel
$W$ of the $L$ block action is contained in the even-weight
space $E$ and is invariant under $S_m$.
If $W$ contains a nonzero constant vector, it contains
$\lambda$. If it contains a nonconstant vector, its difference
with a transposition-conjugate having unequal coordinates
has weight two. Its $S_m$-conjugates generate $E$.
For even $m$, $\lambda\in E$, again contradicting
$\lambda\notin L$. Thus $W=0$ and $L\cong S_m$.
But $H$ fixes $a$ and has order $(m-1)!$, since its block
image is $S_{m-1}$. Transitivity of $L$ on $2m$ points
would instead require
\[
                  |L_a|=\frac{m!}{2m}=\frac{(m-1)!}{2},
\]
which cannot contain $H$. This excludes all even $m$.

For odd $m\geq5$, apply the accepted paired-row centralizer
lemma to $H$ on the other $m-1$ pairs. Their number is even
and at least four, and the block image is $S_{m-1}$.
Every centralizing involution there is even. Since $u$ fixes
$a,a'$, it is even on $A$, making $(a\ b)u$ odd.
This contradicts $G\leq A_n$.

The only remaining value is $m=3$. The two generators of $L$
belong to $E\rtimes S_3$, since they are even on the six
points, and their block images are distinct transpositions
generating $S_3$. Their product has a three-cycle block image.
For any even vector $v\in\mathbb F_2^3$ and a three-cycle
$\pi$, the cube of $(v,\pi)$ has base vector
\[
                   v+\pi v+\pi^2v
                   =\Bigl(\sum_jv_j\Bigr)(1,1,1)=0.
\]
Thus the product of the two involutions has order three.
They generate a group of order at most six; the surjective
block image shows $L\cong S_3$ exactly. Yet $H$ is a
nontrivial order-two subgroup fixing $a$, whereas a transitive
degree-six action of this order-six group would be regular.
This final contradiction proves the proposition.
\end{proof}

\paragraph{Source and boundary.}
This replaces the transitive-side equality-graph step in
Cameron--Fernandes--Leemans--Mixer, \emph{Highest rank of a
polytope for $A_n$}, arXiv:1605.09173v1, Proposition~5.26,
physical page~34. The complete numerical transitive-side
proof, unique-crossing lemma, paired-row parity lemma and
retained Whiston bounds are paid inputs. The binary-kernel
argument is written in full; no finite computation or new
formal theorem is required. This note does not remove the
separate intransitive-side $D_{10}$ residual, prove every
transitive-tail rank bound, or complete the whole L044 source.
No original, question or priority credit is asserted.
\end{document}
